\section{Problem Formulation and Motivation}
\label{sec:formulation}

Consider a dataset
\(
\mathcal{D}=\{(\tau_i,r_i)\}_{i=1}^{N}
\),
where
\(
\tau_i=(q_i,s_{i,1},\ldots,s_{i,T_i})
\)
is a reasoning trajectory and \(r_i\in\{0,1\}\) is its final reward.
We view the steps
\(
\mathcal{N}_i=\{s_{i,1},\ldots,s_{i,T_i}\}
\)
as players in a cooperative game whose payoff is the trajectory outcome.
For a coalition \(S\subseteq\mathcal{N}_i\) with value \(v(S)\), classical
Shapley assigns step \(t\) the expected marginal contribution
\begin{equation}
\label{eq:discrete_shapley}
\phi_{i,t}^{\mathrm{sha}}
:=
\frac{1}{T_i!}
\sum_{\pi\in\Pi(\mathcal{N}_i)}
\Big[
v\!\big(\mathrm{Pre}_{\pi}(s_{i,t})\cup\{s_{i,t}\}\big)
-
v\!\big(\mathrm{Pre}_{\pi}(s_{i,t})\big)
\Big],
\end{equation}
where \(\Pi(\mathcal{N}_i)\) is the set of permutations and
\(\mathrm{Pre}_{\pi}(s_{i,t})\) denotes predecessors of \(s_{i,t}\) under
\(\pi\) \citep{shapley1953value}.

This allocation principle is attractive because it formalizes credit as
marginal contribution.
To apply it to reasoning trajectories, we first ask what Shapley becomes when
the temporal order of the trajectory is respected.
If the observed reasoning order is treated as the admissible causal order, the
allocation reduces to a prefix value difference.
Let
\(
P_{i,t}=\{s_{i,1},\ldots,s_{i,t}\}
\)
and \(P_{i,0}=\emptyset\).
\begin{lemma}
\label{lem:ordered_shapley_degenerate}
If only \(1\prec 2\prec \cdots \prec T_i\) is admissible, then
\Cref{eq:discrete_shapley} reduces to
\(
\phi_{i,t}^{\mathrm{ord}}=v(P_{i,t})-v(P_{i,t-1})
\).
\end{lemma}

The prefix value difference also has a reward-redistribution interpretation.
In an undiscounted episodic setting, let \(v^\star(P_{i,t})\) be the optimal
prefix value and define
\(
\tilde r_{i,t}^\star=v^\star(P_{i,t})-v^\star(P_{i,t-1})
\).
\begin{lemma}
\label{lem:reward_shaping_invariance}
The value-difference reward \(\tilde r_{i,t}^\star\) is a potential-based
redistribution and preserves the optimal policy of the original terminal-reward
problem.
\end{lemma}
\Cref{sec:appendix_local_credit_proofs} provides proof sketches for both
lemmas.

Together, the two lemmas show that Shapley-style marginal contribution is a
valid credit-assignment principle for ordered reasoning.
Under the temporal constraint, Shapley becomes a value-difference reward; under
an ideal prefix value function, that reward can preserve the optimal policy of
the original terminal-reward problem.
This makes Shapley a principled target for step attribution.
However, directly applying this principle still does not remove the fundamental
dependence on an external strong value model or step-level value labels.
We therefore retain Shapley's marginal-contribution rationale while eliminating
the need for externally specified step-level value or reference distributions
\citep{ng1999policy,pmlr-v119-sundararajan20b,ovm:conf/naacl/YuGW24}.
